\section*{Введение}
\addcontentsline{toc}{section}{Введение}

\textbf{Актуальность темы.} В условиях роста интереса к здоровому образу жизни и правильному питанию возрастает потребность в удобных инструментах для контроля рациона. Традиционные методы ведения дневников питания (бумажные блокноты или таблицы) неудобны, требуют значительных временных затрат и не дают оперативной аналитики. Разработка специализированного веб-приложения позволяюет автоматически подсчитывать калорийность блюд, анализировать баланс белков, жиров и углеводов (БЖУ), а также отслеживать динамику питания.

\textbf{Степень разработанности проблемы.} Вопросы анализа питания и расчета энергетической ценности продуктов активно исследуются в области нутрициологии и информационных технологий. В области веб-ориентированных систем существенный вклад внесли создатели популярных приложений (например, FatSecret, Lifesum). Однако многие из них имеют ограничения, связанные с платным доступом, недостаточной адаптацией под индивидуальные потребности пользователей или ограниченной гибкостью. Разработка доступных, расширяемых и адаптивных веб-приложений остается актуальной задачей.

\textbf{Объект исследования} — процессы учета и анализа питания пользователей в рамках цифровых систем.

\textbf{Предмет исследования} — программные средства и алгоритмы, обеспечивающие автоматизацию подсчета калорий, анализ макронутриентов и визуализацию данных в веб-приложении.

\textbf{Целью данной работы} является проектирование и разработка веб-приложения, предназначенного для учета рациона питания пользователей, автоматического расчета калорийности и анализа структуры питания с целью формирования более сбалансированных пищевых привычек.
